% STATUS: written 2026-10-12
% STATUS: extended 2026-11-07 (Awtar 1973, Bresar and Vukman 1993, Golbasi and
% Kaya 2006, Ashraf and Rehman 2000 and Sandhu and Rehman 2025 read in full
% from the library PDFs)
% Sources read in full: Bergen, Herstein and Kerr 1981 (pp. 259--267), Martindale and Miers 1986 (p. 14), Lanski 1988 (pp. 275--276, 280);
% Awtar 1973 (pp. 67--75); Bresar and Vukman 1993 (pp. 133--141); Golbasi and
% Kaya 2006 (pp. 862--866); Ashraf and Rehman 2000 (pp. 201--206); Sandhu and
% Rehman 2025 (pp. 1--12 of the arXiv version).
\section{Derivations acting on ideals and Lie ideals}\label{sec:prime-semiprime:ideals-lie-ideals}

Section~\ref{sec:prime-semiprime:engel-nilpotent} studied a derivation on the
whole ring or on an ideal. We now ask how a derivation acts on a Lie ideal. A
\emph{Lie ideal} of $R$ is an additive subgroup $U$ with $[U,R]\subseteq U$
\cite[p.~280]{Lanski1988}.

Herstein began to study the Jordan and Lie
ideals of a simple ring in the early 1950s \cite[p.~14]{MartindaleMiers1986}. The
work was then carried to prime and semiprime rings, and Lanski was a major force
in it \cite[p.~14]{MartindaleMiers1986}.

Bergen, Herstein and Kerr took up derivations and Lie ideals together
\cite[p.~259]{BergenHersteinKerr1981}. They measure the size of $d(U)$ in two
ways \cite[p.~259]{BergenHersteinKerr1981}. The first way is the centralizer of $d(U)$: the bigger $d(U)$ is, the smaller its
centralizer should be \cite[p.~259]{BergenHersteinKerr1981}. The other way asks
whether the subring generated by $d(U)$ contains a nonzero ideal of $R$
\cite[p.~259]{BergenHersteinKerr1981}.

They work in a prime ring $R$ of characteristic not $2$, with center $Z$
\cite[p.~260]{BergenHersteinKerr1981}. Their first tool turns a Lie ideal into an
ideal.

\begin{theorem}[{\cite[Lemma~1, p.~260]{BergenHersteinKerr1981}}]\label{thm:prime:bhk-lemma}
Let $R$ be a prime ring of characteristic not $2$ and let $U\not\subseteq Z$ be a
Lie ideal of $R$. Then there is an ideal $M$ of $R$ with $[M,R]\subseteq U$ and
$[M,R]\not\subseteq Z$.
\end{theorem}

We state the first theorem, on the case $d^2(U)=0$.

\begin{theorem}[{\cite[Theorem~1, p.~261]{BergenHersteinKerr1981}}]\label{thm:prime:bhk-one}
Let $R$ be a prime ring of characteristic not $2$, let $d\neq 0$ be a derivation
of $R$, and let $U$ be a Lie ideal of $R$ with $d^2(U)=0$. Then $U\subseteq Z$.
\end{theorem}

We state the second theorem, on the centralizer of $d(U)$.

\begin{theorem}[{\cite[Theorem~2, p.~262]{BergenHersteinKerr1981}}]\label{thm:prime:bhk-two}
Let $R$ be a prime ring of characteristic not $2$, let $U\not\subseteq Z$ be a Lie
ideal of $R$, and let $d\neq 0$ be a derivation of $R$. Then the centralizer of
$d(U)$ in $R$ is $Z$.
\end{theorem}

We state the third theorem, on the size of $d(U)$.
\begin{theorem}[{\cite[Theorem~3, p.~265]{BergenHersteinKerr1981}}]\label{thm:prime:bhk-three}
Let $R$ be a prime ring of characteristic not $2$, let $U\not\subseteq Z$ be a Lie
ideal of $R$, and let $d$ be a derivation of $R$ with $d^3\neq 0$. Then the
subring generated by $d(U)$ contains a nonzero ideal of $R$.
\end{theorem}

We state the last theorem, on the product of two derivations.

\begin{theorem}[{\cite[Theorem~4, p.~266]{BergenHersteinKerr1981}}]\label{thm:prime:bhk-four}
Let $R$ be a prime ring of characteristic not $2$ and let $U\not\subseteq Z$ be a
Lie ideal of $R$. Let $\delta$ and $d$ be derivations of $R$ with
$\delta d(U)=0$. Then $d=0$ or $\delta=0$.
\end{theorem}

Theorem~\ref{thm:prime:bhk-one} is a special case of
Theorem~\ref{thm:prime:bhk-four} with $\delta=d$
\cite[p.~266]{BergenHersteinKerr1981}. Several authors proved that $d=0$ or $h=0$
when the product of two derivations vanishes on a Lie ideal or lies in the center
\cite[p.~275]{Lanski1988}.

Lanski removed the restriction on the characteristic
\cite[p.~276]{Lanski1988}. He treats a noncommutative Lie ideal in every
characteristic. He notes one exception, a ring of characteristic $2$ that satisfies the
standard identity of degree four. Section~\ref{sec:prime-semiprime:posner}
stated these results.

Awtar had already put the hypothesis of the second theorem of Posner on a Lie ideal
and on a Jordan ideal \cite[p.~67]{Awtar1973}. He takes a nonzero derivation $d$ of
a prime ring and asks that $[u,d(u)]$ lie in the center for every $u$ in the ideal
\cite[p.~68]{Awtar1973}. A Jordan ideal is an additive subgroup $U$ with
$u\circ r\in U$ for all $u\in U$ and $r\in R$, where $u\circ r=ur+ru$
\cite[p.~71]{Awtar1973}.

\begin{theorem}[{\cite[Theorem~1, p.~68]{Awtar1973}}]\label{thm:prime:awtar-lie}
Let $R$ be a prime ring of characteristic different from $2$ and $3$, let $d$ be
a nonzero derivation of $R$, and let $U$ be a Lie ideal of $R$ with $[u,d(u)]\in Z$
for all $u\in U$. Then $U\subseteq Z$.
\end{theorem}

Awtar notes that the exclusion of characteristic $3$ does not enter the proof after
the first step \cite[p.~71]{Awtar1973}. So he states the characteristic $3$ case
separately for a square closed Lie ideal, one with $u^{2}\in U$ for all $u\in U$.

\begin{theorem}[{\cite[Theorem~2, p.~71]{Awtar1973}}]\label{thm:prime:awtar-char3}
Let $R$ be a prime ring of characteristic $3$ and $d$ a nonzero derivation of $R$.
If $U$ is a Lie ideal of $R$ with $[u,d(u)]\in Z$ and $u^{2}\in U$ for all $u\in U$,
then $U\subseteq Z$.
\end{theorem}

The same conclusion holds for a Jordan ideal away from characteristic $2$
\cite[p.~71]{Awtar1973}.

\begin{theorem}[{\cite[Theorem~3, p.~71]{Awtar1973}}]\label{thm:prime:awtar-jordan}
Let $R$ be a prime ring of characteristic not $2$, let $d$ be a nonzero derivation
of $R$, and let $U$ be a Jordan ideal of $R$ with $[u,d(u)]\in Z$ for all $u\in U$.
Then $U\subseteq Z$.
\end{theorem}

In characteristic $2$ the Lie and Jordan ideals coincide, and Awtar proves a weaker
statement for a Lie ideal that is also a subring \cite[p.~72]{Awtar1973}.

\begin{theorem}[{\cite[Theorem~4, p.~72]{Awtar1973}}]\label{thm:prime:awtar-char2}
Let $R$ be a prime ring of characteristic $2$, let $d$ be a nonzero derivation of
$R$, and let $U$ be a Lie ideal and a subring of $R$ with $[u,d(u)]\in Z$ for all
$u\in U$. Then $U$ is commutative.
\end{theorem}

The condition on $Z$ describes the size of the set of values. Bre\v{s}ar and Vukman
measured that set directly \cite[p.~133]{BresarVukman1993}. They take the subring
generated by all the elements $[D(x),x]$, and they prove that it contains a
one-sided ideal \cite[p.~133]{BresarVukman1993}.

\begin{theorem}[{\cite[p.~133]{BresarVukman1993}}]\label{thm:prime:bv-subring}
Let $R$ be a noncommutative prime ring of characteristic not $2$, and let $D$ be a
nonzero derivation of $R$. Then $U$, the subring of $R$ generated by all
$[D(x),x]$ for $x\in R$, contains a nonzero left ideal of $R$ and a nonzero right
ideal of $R$.
\end{theorem}

A noncommutative prime ring has no nonzero central one-sided ideal, so the theorem
strengthens the statement of Posner that $U$ is not contained in the center
\cite[p.~133]{BresarVukman1993}. The proof rests on a lemma about two derivations.

\begin{lemma}[{\cite[Lemma~1, p.~134]{BresarVukman1993}}]\label{lem:prime:bv-two}
Let $R$ be a prime ring of characteristic not $2$. If there exist nonzero
derivations $D$ and $G$ of $R$ such that $G([D(x),x])=0$ for all $x\in R$, then $R$
is commutative.
\end{lemma}

The lemma fails in characteristic $2$, and the authors give an example with an
inner derivation \cite[p.~134]{BresarVukman1993}. They close with a question: does
$U$ contain a nonzero two-sided ideal \cite[p.~141]{BresarVukman1993}?

G\"olba\c{s}{\i} and Kaya carried the same questions to a generalized derivation
\cite[p.~862]{GolbasiKaya2006}. On a prime ring of characteristic not $2$ they show
that the values control the Lie ideal in three ways
\cite[Theorem~1, p.~864]{GolbasiKaya2006}. The first two results are the following.

\begin{theorem}[{\cite[Theorem~1, p.~864]{GolbasiKaya2006}}]\label{thm:prime:gk-square}
Let $R$ be a prime ring of characteristic not $2$, let $U$ be a nonzero Lie ideal
of $R$, and let $f$ be a generalized derivation of $R$ associated with a derivation
$d$. If $f^{2}(U)=0$, then $U\subseteq Z$.
\end{theorem}

Their lemma that a nonzero $f$ cannot vanish on a noncentral Lie ideal stands
before it \cite[Lemma~5, p.~863]{GolbasiKaya2006}. They also treat a map that acts
as a homomorphism on the Lie ideal, under the extra condition $u^{2}\in U$ for all
$u\in U$ \cite[p.~864]{GolbasiKaya2006}.

\begin{theorem}[{\cite[Theorem~2, p.~864]{GolbasiKaya2006}}]\label{thm:prime:gk-hom}
Let $R$ be a prime ring of characteristic not $2$, let $U$ be a Lie ideal of $R$
with $u^{2}\in U$ for all $u\in U$, and let $f$ be a generalized derivation of $R$
associated with a derivation $d$. If $f(uv)=f(u)f(v)$ for all $u,v\in U$, then
$U\subseteq Z$.
\end{theorem}

The same conclusion holds when $f$ acts as an antihomomorphism on $U$, and when
$f(uv)=f(vu)$ for all $u,v\in U$
\cite[Theorems~3 and~4, p.~864]{GolbasiKaya2006}.

A Jordan left derivation is the map of
Definition~\ref{def:gen:left-derivation} with $D(u^{2})=2uD(u)$. Ashraf and
Rehman put that identity on a square closed Lie ideal
\cite[p.~200]{AshrafRehman2000}. Their map $d$ is additive on the whole ring and
the identity is assumed only on the Lie ideal, not on all of $R$
\cite[p.~205]{AshrafRehman2000}.

\begin{theorem}[{\cite[p.~205]{AshrafRehman2000}}]\label{thm:prime:ashraf-jordan-left}
Let $R$ be a $2$-torsion free prime ring and let $U$ be a Lie ideal of $R$ with
$u^{2}\in U$ for all $u\in U$. If $d\colon R\to R$ is an additive map with
$d(u^{2})=2ud(u)$ for all $u\in U$, then $d(uv)=ud(v)+vd(u)$ for all $u,v\in U$.
\end{theorem}

So on such a Lie ideal a Jordan left derivation obeys the full left rule. The
proof uses the square closed condition at every step, since the sums and products
it forms must stay in $U$ \cite[p.~205]{AshrafRehman2000}.

Sandhu and Rehman replace the prime ring by a ring with a semiprime ideal and put
the central values into that ideal \cite[p.~1]{SandhuRehman2025}. Let $T$ be a
semiprime ideal of a ring $R$ with $T\subsetneq I$ for a nonzero ideal $I$. A
derivation $d$ is \emph{$T$-commuting} on a set $K$ when $[d(x),x]\in T$ for all
$x\in K$ \cite[p.~1]{SandhuRehman2025}.

\begin{theorem}[{\cite[Theorem~2.1, p.~2]{SandhuRehman2025}}]\label{thm:prime:sandhu-engel}
Let $R$ be a ring with a nonzero ideal $I$ and a semiprime ideal $T$ such that
$T\subsetneq I$. If $\operatorname{char}(R/T)\neq 2$ and $[[d(x),x],x]\in T$ for
all $x\in I$, then $d$ is $T$-commuting on $I$.
\end{theorem}

The authors note that their argument is shorter than the one of Vukman for the same
identity on a prime ring \cite[p.~2]{SandhuRehman2025}. Two corollaries recover the
known statements over a prime ring \cite[p.~4]{SandhuRehman2025}. The second of
them concerns a noncommutative prime ring of characteristic not $2$ on which
$x\mapsto[d(x),x]$ is centralizing. There $d=0$, and the exclusion of
characteristic $3$ from a theorem of Vukman is dropped
\cite[Corollary~2.7, p.~6]{SandhuRehman2025}. Their next theorem puts the values in
the center of the quotient $R/T$ instead of in $T$, and reaches the same conclusion
\cite[Theorem~2.5, p.~5]{SandhuRehman2025}.

Section~\ref{sec:prime-semiprime:involution} asks the same question for the skew
elements of a ring with involution, which form a Lie subring.

A Lie derivation can be tested on the products that vanish. Qi, Cui and Hou
characterize the additive $\xi$-Lie derivations of a prime algebra by their
action on $\xi$-Lie zero products \cite[pp.~669--682]{QiCuiHou2011XiLie}.

